% ===========================================================================
\section{Introdu\c{c}\~ao}
\label{sec:introducao}
% ===========================================================================

A modelagem tridimensional \'e, na engenharia contempor\^anea, uma atividade de projeto e n\~ao
apenas de ilustra\c{c}\~ao. Um modelo de superf\'icie de ve\'iculo automotivo codifica decis\~oes
geom\'etricas --- curvaturas, tangentes, dimens\~oes de envelope, folgas --- que podem ser
medidas, verificadas e reproduzidas. Quando essa geometria nasce de um fluxo manual e
interativo, por\'em, boa parte dessas decis\~oes permanece impl\'icita: vive na mem\'oria de
quem operou o \emph{software}, e n\~ao em um artefato audit\'avel. Um ajuste de malha
bem-sucedido pode n\~ao ser reconstru\'ido depois, porque foi feito com o mouse, em uma
sequ\^encia de gestos que ningu\'em registrou \cite{bbotsch2010}.

Em paralelo, os protocolos de integra\c{c}\~ao entre modelos de linguagem de grande porte e
ferramentas externas amadureceram rapidamente. O \emph{Model Context Protocol} (MCP)
\cite{mcpspec2024} prop\~oe uma interface padronizada, baseada em mensagens JSON-RPC~2.0
\cite{jsonrpc2020}, para que um agente descubra e invoque ferramentas expostas por um
servidor. Um servidor MCP para o Blender \cite{blendermcp2025} permite que um agente execute
c\'odigo Python na sess\~ao viva do programa, inspecione a cena e capture imagens da
\emph{viewport}. A unidade de trabalho deixa de ser o gesto e passa a ser o
\emph{script}, que \'e version\'avel, revis\'avel e reproduz\'ivel.

A conflu\^encia desses dois movimentos --- formaliza\c{c}\~ao da geometria e automa\c{c}\~ao por
agentes --- permite tratar a modelagem de um ve\'iculo como um problema de engenharia de
\emph{software}: com requisitos identificados, decis\~oes de arquitetura registradas,
crit\'erios de aceita\c{c}\~ao observ\'aveis e evid\^encia de verifica\c{c}\~ao rastre\'avel
\cite{shah1995, piegl1997, pahl2007}.

\subsection{Problema e lacuna}

O problema que orienta este trabalho pode ser enunciado assim:

\begin{quote}
\emph{\'E poss\'ivel produzir, de forma integralmente reproduz\'ivel e automatizada por um
agente de intelig\^encia artificial por meio do Model Context Protocol, um modelo
tridimensional de um autom\'ovel esportivo cuja geometria e materiais sejam fi\'eis a
refer\^encias fotogr\'aficas dispon\'iveis, atendendo simultaneamente a restri\c{c}\~oes de
topologia de malha adequada \`a subdivis\~ao e a um envelope dimensional verificado?}
\end{quote}

Do enunciado derivam tr\^es tens\~oes t\'ecnicas que precisam ser resolvidas conjuntamente. A
primeira \'e a \textbf{tens\~ao entre fidelidade e automa\c{c}\~ao}: refer\^encias
fotogr\'aficas s\~ao perspectivas, com distor\c{c}\~ao e oclus\~ao, e n\~ao \emph{blueprints}
ortogr\'aficos; converter fotografias em par\^ametros geom\'etricos exige um m\'etodo
expl\'icito de estimativa \cite{hartley2004}. A segunda \'e a \textbf{tens\~ao entre
expressividade e restri\c{c}\~ao topol\'ogica}: reproduzir recortes, arcos de roda e
entradas de ar costuma ser feito com opera\c{c}\~oes booleanas, que produzem n-gons e
tri\^angulos exatamente sobre arestas vis\'iveis, incompat\'iveis com subdivis\~ao de
qualidade \cite{catmull1978, bbotsch2010}. A terceira \'e a \textbf{tens\~ao entre realismo
f\'isico e custo computacional}: pintura automotiva convincente requer
\emph{clear coat}, micro-flocos met\'alicos e transmiss\~ao em vidros, o que demanda um
modelo de sombreamento fisicamente plaus\'ivel \cite{burley2012, karis2013, walter2007}.

A lacuna que este artigo pretende ocupar \'e de natureza metodol\'ogica. A literatura
descreve, de um lado, gera\c{c}\~ao de forma por regras e, de outro, agentes capazes de
invocar ferramentas; o que ainda \'e escasso s\~ao \emph{relatos de caso com
verifica\c{c}\~ao externa, nos quais se registre o que a media\c{c}\~ao por agente
efetivamente agrega --- e, igualmente importante, o que ela n\~ao agrega}.

\subsection{Hip\'otese e justificativa}

A hip\'otese de trabalho \'e que a combina\c{c}\~ao de tr\^es elementos --- (i) uma
especifica\c{c}\~ao formal com requisitos e crit\'erios de aceita\c{c}\~ao observ\'aveis,
(ii) gera\c{c}\~ao de geometria exclusivamente por c\'odigo determin\'istico, com as
restri\c{c}\~oes topol\'ogicas impostas no n\'ivel do algoritmo, e (iii) media\c{c}\~ao de
todas as opera\c{c}\~oes destrutivas por um agente conectado via MCP --- permite obter um
modelo cuja ader\^encia \`a refer\^encia \'e compar\'avel \`a de um fluxo manual equivalente,
mas cuja reprodutibilidade e auditabilidade s\~ao substancialmente superiores.

A justificativa apoia-se em tr\^es argumentos. \textbf{Reprodutibilidade}: um modelo gerado
por \emph{script} \'e reconstru\'ido a partir do reposit\'orio com um \'unico comando,
eliminando a depend\^encia de sess\~oes manuais \cite{shah1995}. \textbf{Auditabilidade do
processo de IA}: quando um agente opera uma ferramenta, o protocolo define exatamente o que
foi pedido e o que foi executado, criando um registro independente do fornecedor do modelo
de linguagem \cite{schick2023, patil2023, yao2023}. \textbf{Forma\c{c}\~ao} de engenharia:
derivar requisitos de um documento informal, convert\^e-los em crit\'erios observ\'aveis e
comprovar o atendimento por evid\^encia reproduz\'ivel \'e a compet\^encia central do projeto
de engenharia \cite{pahl2007, otto2001}.

\subsection{Objetivo, contribui\c{c}\~oes e delimita\c{c}\~ao}

O objetivo \'e produzir um modelo tridimensional do Lotus Elise 111S (Series~1) em escala
real, constru\'ido inteiramente por c\'odigo executado em Blender por meio de um agente
operando pelo MCP, com fidelidade morfol\'ogica, integridade topol\'ogica e plausibilidade
f\'isica dos materiais verificadas por evid\^encia. As contribui\c{c}\~oes reivindicadas
s\~ao quatro: um gerador procedimental completo e reproduz\'ivel de um ve\'iculo
esportivo, com topologia controlada por constru\c{c}\~ao; uma t\'ecnica de tampa de
\emph{loft} que emite apenas quads; uma demonstra\c{c}\~ao de arquitetura de agentes mediada
por MCP para conte\'udo tridimensional, com registro do que ela agrega e do que n\~ao
agrega; e um conjunto de evid\^encia visual comparativa com resultado registrado por
crit\'erio, \emph{incluindo os desvios}.

A pesquisa delimita-se a um \'unico ve\'iculo, uma \'unica variante e um \'unico
acabamento. Est\~ao fora de escopo, por decis\~ao expl\'icita: su\'ite de testes
automatizados; valida\c{c}\~ao de continuidade por mapas de zebra e isofotas; o \emph{rig}
de ilumina\c{c}\~ao de est\'udio completo; interior detalhado; decais e
anima\c{c}\~ao; e exporta\c{c}\~ao para formatos de interc\^ambio. A
verifica\c{c}\~ao \'e visual e comparativa, com registro de aceite por crit\'erio ---
escolha coerente com a natureza do artefato, j\'a que a qualidade de uma superf\'icie
automotiva \'e julgada pela continuidade dos reflexos \cite{halstead1993}. O artigo segue
com fundamenta\c{c}\~ao (Se\c{c}\~ao~\ref{sec:fundamentacao}), trabalhos relacionados
(Se\c{c}\~ao~\ref{sec:relacionados}), materiais e m\'etodos
(Se\c{c}\~ao~\ref{sec:metodos}), desenvolvimento
(Se\c{c}\~ao~\ref{sec:desenvolvimento}), resultados
(Se\c{c}\~ao~\ref{sec:resultados}), discuss\~ao (Se\c{c}\~ao~\ref{sec:discussao}) e
conclus\~ao (Se\c{c}\~ao~\ref{sec:conclusao}).
